% !TeX root = ../prompt-engineering-guide-for-beginners.tex

%%%%%%%%%%%%%%%%%%%%%%%%%%%%%%%%%%
%% 第四章：提示词要素
%%%%%%%%%%%%%%%%%%%%%%%%%%%%%%%%%%
\chapter{提示词要素：一条好 Prompt 的解剖课}
\label{ch:prompt-elements}

了解了 AI 的基本工作原理之后，我们来聊聊提示词到底应该「长什么样」。
一条好的提示词并不是随手一写那么简单，它有自己内在的结构和要素。
搞清楚这些要素，你就掌握了写好提示词的底层方法论。

\section{提示词的四大要素}

一条完整的提示词，最多可包含四个核心组成部分。虽然不是每次对话都要全部写满，但了解它们能帮你在需要时，像搭积木一样快速组装出高质量的指令。

\subsection*{指令（Instruction）}

这是提示词的核心：明确告诉 AI 你想让它做什么。

好的指令要用具体的动词，避免笼统含糊。比如：

\begin{itemize}
  \item 「总结」而不是「说说」
  \item 「列出 5 个要点」而不是「随便讲讲」
  \item 「翻译成英文」而不是「帮我弄一下」
\end{itemize}

\subsection*{上下文（Context）}

背景信息能够帮助 AI 更好地理解你的需求。你提供的上下文越充分，AI 的回答就越贴合你的实际应用场景。

比如同样是「帮我写一封邮件」，在这基础上加上「这是发给甲方的催款邮件，我们合作了三年，关系一直不错，但这次对方拖了两个月没付款」，AI 就能极其精准地拿捏住「既不撕破脸又要施加压力」的妥帖语气。

\subsection*{输入数据（Input Data / Payload）}

如果你需要 AI 处理一段具体的内容，就需要把这段内容干净地注入提示词中。比如：

\begin{itemize}
  \item 需要处理提取知识点的数百行源文件代码
  \item 需要总结的文章文本（Context Corpus）
  \item 需要做可视化推理的脏数据表（CSV）
  \item 需要改写的初稿段落
\end{itemize}

\begin{tipbox}
从工程化角度来看，Prompt 就是一段由“操作代码 (Instruction)”和“数据参数 (Payload)”组成的代码块。为了防止 AI 在读取时把你的文章内容误认为是新下达的指令（即常见的 Prompt Injection 提示词注入漏洞），强烈建议使用特定的 XML 标签（如 \verb|<text>...</text>|）或三个反引号（\verb|```|）将输入数据严格包裹隔离起来。
\end{tipbox}

\subsection*{输出格式（Output Format / Schema Constraints）}

告诉 AI 你期望回答以何种结构形式呈现。作为严谨的提示词工程师，你应当将其理解为与 API 协同中的接口 Schema 定义：

\begin{itemize}
  \item 「以 Markdown 表格形式展示」
  \item 「用层级清晰的编号列表列出」
  \item 「仅仅输出合法的 JSON 格式字符串，不需要任何解释」
  \item 「使用 LaTeX 公式语法包裹所有数学计算过程」
  \item 「强硬控制输出不超过 300 个 Token」
\end{itemize}

\section{万能公式：角色 + 背景 + 任务 + 格式要求}
\label{sec:prompt-formula}

如果你觉得四大要素不太好记，这里有一个更简单的万能公式。写提示词的时候，按照这个顺序填空就行：

\begin{enumerate}
  \item \textbf{角色}：让 AI 扮演什么身份？
  \item \textbf{背景}：这件事的背景是什么？
  \item \textbf{任务}：你需要 AI 做什么？
  \item \textbf{格式要求}：你希望回答是什么样子的？
\end{enumerate}

来看一个实际例子。

\textbf{普通版：}

\begin{promptbox}
帮我写个请假条。
\end{promptbox}

AI 可能会写出一个非常模板化、没有细节的请假条。

\textbf{万能公式版：}

\begin{promptbox}
你是一名 500 强企业的员工（角色），由于重感冒需要向严厉的老板请假两天（背景），请帮我写一封请假邮件（任务），要求态度诚恳、措辞专业、字数 100 字以内（格式要求）。
\end{promptbox}

用了万能公式之后，AI 会给出一封语气恰当、内容具体、长度合适的请假邮件。

不需要每次都写得这么面面俱到，但当你对 AI 的回答感到不满时，可以回过头来拿着这个公式检查一下：是不是缺了角色？是不是没交代背景？还是没说清楚输出的格式？

再来一个接地气的例子：

\begin{promptbox}
你是一位上海本地美食达人（角色），我外地朋友下周来上海玩三天（背景），请推荐五家不踩雷的本帮菜馆，涵盖从人均 50 到人均 300 的不同档次（任务），用表格列出餐厅名、地址、人均消费和推荐菜品（格式要求）。
\end{promptbox}

这样一写，AI 就不会给你推荐一堆连锁快餐或者搜索结果里千篇一律的「十大必吃」了，而是会按照你要求的表格格式、价位梯度，给出更有参考价值的推荐。

\section{好提示词 vs 坏提示词：对比拆解}

下面通过几组对比，让你直观感受提示词质量对结果的影响。

\subsection*{案例一：写一篇产品介绍}

\textbf{坏提示词：}
\begin{promptbox}
帮我写个产品介绍。
\end{promptbox}

\textbf{问题}：AI 不知道是什么产品，也不知道写给谁看，只能写出一篇空洞的通用模板。

\textbf{好提示词：}
\begin{promptbox}
请帮我写一段产品介绍。产品是一款面向年轻白领的便携式咖啡机，主打卖点是 30 秒快速萃取和超轻便携设计。介绍文案发布在小红书上，语气活泼种草，控制在 150 字以内。
\end{promptbox}

\textbf{差别}：有了具体产品、目标人群、发布渠道、语气和字数要求，AI 就能精准输出。

\subsection*{案例二：翻译一段文字}

\textbf{坏提示词：}
\begin{promptbox}
翻译这个：We need to leverage our core competencies to drive synergistic outcomes.
\end{promptbox}

\textbf{好提示词：}
\begin{promptbox}
请把下面这句英文翻译成中文，保持商务风格，避免直译，用自然流畅的表达：

We need to leverage our core competencies to drive synergistic outcomes.
\end{promptbox}

\textbf{差别}：加上翻译风格和要求，AI 的翻译会更加自然地道，而不是逐字硬译。

\subsection*{案例三：做一个工作计划}

\textbf{坏提示词：}
\begin{promptbox}
帮我做个计划。
\end{promptbox}

\textbf{好提示词：}
\begin{promptbox}
我是一名新媒体运营，下个月需要完成公司公众号的日常运营。请帮我制定一份为期 4 周的内容发布计划，要求每周发布 3 篇文章，涵盖行业资讯、产品介绍和用户故事三个主题，以表格形式呈现，包含日期、主题、标题建议和关键词。
\end{promptbox}

\textbf{差别}：拥有了详细的背景交代、清晰的任务目标和明确的表格输出格式，AI 便能直接产出一份可落地执行的高质量计划。

通过这几组对比可以看出，写好提示词的核心其实就是三个字：\textbf{说清楚}。你给 AI 的信息越清晰、越具体，它的输出就越贴近你的预期。

\section{快速检查清单}

每次写完提示词、准备发送之前，可以用下面这个清单快速自查一遍：

\begin{center}
\begin{tabular}{lp{9cm}}
  \toprule
  \textbf{检查项} & \textbf{问自己} \\
  \midrule
  任务明确 & 我是否用了具体的动词（总结/翻译/列出），而不是「帮我看看」？ \\
  背景充分 & AI 是否知道这件事的前因后果和使用场景？ \\
  受众清晰 & AI 是否知道这个输出是给谁看的？ \\
  格式要求 & 我是否说明了输出的形式（表格/列表/字数）？ \\
  约束到位 & 我是否给出了语气、长度、风格等限制条件？ \\
  \bottomrule
\end{tabular}
\end{center}

不必每次都严格检查所有项目，简单的问题一两句话就够了。但当 AI 的回答不尽如人意时，回来对照这个清单，通常很快就能找到原因。

现在你已经知道了一条好提示词应该包含哪些「结构要素」，接下来我们将进一步聊聊该「如何表达」。下一章会为你介绍一系列实用的通用写作技巧。
